\paragraph{Depth-damage functions}
Depth-damage functions map a hazard intensity measure to a damage ratio (repair cost relative to replacement value) and are the primary tool for translating flood hazard into economic loss. In the flood context, inundation depth is the dominant intensity measure and serves as the primary input to these functions~\cite{merz2010review, huizinga2017global}; other intensity measures such as flow velocity and inundation duration are secondary but can be material for specific flood types and asset classes~\cite{merz2010review}. The most widely used functions originate from engineering studies by the US Army Corps of Engineers~\cite{usace2003}, which developed stage-damage curves for residential structures based on expert judgement and post-disaster surveys. The Joint Research Centre of the European Commission~\cite{huizinga2017global} extended this approach globally, producing continent-specific depth-damage curves for residential, commercial, and industrial assets. These JRC functions are the default vulnerability module in CLIMADA~\cite{aznar2019climada}, the open-source platform developed at ETH Zurich for probabilistic natural catastrophe impact assessment.

A common limitation of these functions is the absence of a duration parameter. All standard depth-damage curves treat damage as an instantaneous mapping from water depth to loss, implicitly assuming that a flood of one metre depth produces the same damage whether it recedes in two hours or persists for two weeks. This assumption is physically unrealistic: prolonged inundation causes additional degradation through water absorption by building materials, mould colonisation, corrosion of electrical and mechanical systems, and progressive structural weakening~\cite{merz2010review}. Several engineering studies have documented the duration effect qualitatively and incorporated it into multi-variable damage models alongside depth and flow velocity~\cite{merz2010review}. Recent work has begun to formalise depth-damage relationships for the built environment~\cite{pistrika2014flood} and to extend them toward probabilistic vulnerability and functionality-loss frameworks~\cite{pandey2026urban, forte2025flash, gautam2022unzipping}. Post-event field surveys further document the breadth of damage mechanisms activated during extreme flood episodes~\cite{santangelo2021landslides, thapa2020catchment, santo2016post, chinh2016mekong}. However, quantitative calibration of a duration-dependent parametric function on large-scale insurance claims data has been lacking. This is the gap addressed by the present paper.

\paragraph{Calibration data gap}
The Basel Committee~\cite{basel2021a, basel2021b} and industry surveys~\cite{unepfi2025bridging} consistently identify data scarcity as the primary obstacle to developing reliable physical risk models. Damage function calibration requires paired observations of hazard intensity and economic loss at the building or asset level --- data that are rarely available outside the insurance sector. The FEMA National Flood Insurance Program (NFIP) represents one of the few large-scale sources of such data, with millions of individual claims containing water depth, payout, and building characteristics. Machine learning approaches have demonstrated the potential for multi-variable flood damage estimation when richer hazard and building-characteristic data are available, including gradient-boosted and causal ML models applied to residential properties~\cite{museru2025improving, museru2024advancing} and ensemble deep learning models for flood forecasting~\cite{kordani2024improving}. To our knowledge, this paper presents the first systematic calibration of a parametric damage function with duration dependence on NFIP data at the scale of 295{,}637 claims.

\paragraph{Flood risk across infrastructure sectors}
Flood risk assessment extends beyond residential properties to critical infrastructure and community-scale systems. Studies of wastewater treatment facilities~\cite{nazari2022hydrodynamic, sun2021comprehensive}, coastal structures~\cite{fahad2020coupled, fahad2022decision}, and earthen infrastructure~\cite{motamedi2018quantitative} document sector-specific vulnerability patterns. Economic resilience analyses at the community scale~\cite{nazari2025small} complement asset-level damage estimation by capturing wider socioeconomic disruption. These application contexts motivate the sector adjustment module ($\kappa$) introduced in Section~\ref{ssec: sector}.

\paragraph{Physical climate risk and credit}
A growing body of research demonstrates that physical climate hazards materially affect corporate creditworthiness. Cathcart et al.~\cite{cathcart2026rain} find that temperature and precipitation anomalies increase default probabilities for small firms by 32--87 basis points per standard deviation. Abinzano et al.~\cite{abinzano2026physical} provide global evidence that physical climate risk increases banks' credit risk, with effects concentrated in weakly capitalised institutions. The Bank for International Settlements~\cite{bis2025physical} argues that physical risks have evolved from a niche concern into a systemic risk factor requiring integration into credit risk models.

On the methodology side, Moody's Analytics~\cite{bocchio2023quantifying} and MSCI~\cite{msci2024transition} have developed proprietary frameworks that translate climate scenarios into changes in expected default frequency, using top-down GDP allocation and structural credit models. Bourgey et al.~\cite{bourgey2025climacred} contribute a micro-founded theoretical model linking default probabilities to NGFS scenarios. However, these approaches rely on aggregate damage estimates rather than empirically calibrated, asset-level damage functions --- precisely the gap that our calibration addresses.

\paragraph{Sector-specific vulnerability}
The distinction between residential and commercial/industrial vulnerability is well documented in the flood damage literature~\cite{merz2010review, huizinga2017global}. Commercial buildings typically have higher floor elevations, more robust construction, and operational redundancy that reduces effective damage ratios. However, quantitative estimates of the residential-to-commercial damage ratio---what we term the sector adjustment multiplier $\kappa$---are sparse and typically based on engineering judgement rather than loss data. Our approach calibrates $\kappa$ empirically from 82{,}650 NFIP non-residential claims grouped by observable building characteristics, providing a data-driven alternative to expert-based estimates.
